% Acknowledgements, with the disclosure of AI use arXiv asks authors to make
% (written 2026-10-15 at the owner's request). Software versions from the
% repository: lean/lean-toolchain, the Python environment, Appendix D.
\section*{Acknowledgements}
\addcontentsline{toc}{section}{Acknowledgements}

\paragraph{Use of AI tools.} Much of this work was carried out with Claude Code
(Anthropic), an AI coding assistant running Claude models. Under the authors'
direction it wrote most of the C and Python software, the Lean proofs and drafts of the
text, ran the computations, and also searched and summarised the literature. The
authors chose the questions, directed the work, read the results and take full
responsibility for the content of this paper. We did not take the AI's word for any result: every theorem the paper states
is checked by the Lean kernel (unless the text says otherwise); every value in
the dataset carries a witness that an independent checker verifies; and the
refutations behind the central claims (the wrong optimum of
Theorem~\ref{thm:wrong}, and the minor bounds for the Mycielski graphs and four
MOSP instances) are re-checked by methods independent of our search: a SAT solver with a DRAT
proof that drat-trim verifies, and a separate exact search with its own
checker, written from scratch.
The other refutations rest on our corrected search, whose rules are proved
sound in Lean (Sections \ref{sec:search} and \ref{sec:results}).

\paragraph{Software.} The formal proofs use Lean~4 \citep{MouraUllrich2021},
version 4.34, with the Mathlib library \citep{Mathlib2020}. The SAT checks use
Kissat~4.0.4 (\url{https://github.com/arminbiere/kissat}) and
CaDiCaL~1.9.5 \citep{CaDiCaL2024} through PySAT \citep{PySAT2018}, and the
proofs of unsatisfiability are checked by drat-trim \citep{WetzlerHeuleHunt2014}.
Isomorphism classes are computed with nauty \citep{McKayPiperno2014}, through
pynauty, and the graphs are handled with NetworkX. The bibliometric counts use
OpenAlex \citep{Priem2022}. SageMath provides the openly available
implementation of the branch and bound of \citet{CoudertMazauricNisse2016}.
